%% toc: P3 38 38.7
\chapter{Differences from Ordinary First-Order and Probabilistic Semantics}
\label{sec:p03-w03-p03-053-differences}

Several logics in this part attach a topology, a measure or a probability
to a semantic definition. They do not attach it to the same object.
In one regime a topology lives on the carrier of a single structure and
formulas quantify over its open sets. In another a topology lives on the
space of codes of all countable structures, and no formula mentions it.
In a third a measure or a Baire category notion lives on the carrier of a
single structure and a new quantifier says that a definable section is
large. In a fourth a probability lives on the set of all labelled finite
structures of one size, and the formulas are ordinary first-order
sentences. This chapter states, for each regime, the logical language,
the class of structures, the semantic parameter and where that parameter
lives. It then proves which ordinary first-order results survive, which
change form, and which fail.

\emph{Inputs.} Ordinary first-order logic and its compactness,
completeness, downward L\"owenheim--Skolem and interpolation theorems
(\cref{thm:fol-compactness,thm:fol-completeness,%
thm:downward-ls,thm:p02-craig-interpolation});
\(L_{\omega_1\omega}\) (\cref{def:p02-lomega-one-omega});
the topological languages \(L_2\) and \(L_t\)
(\cref{def:p03-topological-carrier,def:p03-topological-languages});
the coding space \(X_L\) (\cref{def:p03-morley-coding});
the zero--one law (\cref{thm:p03-fol-zero-one});
and the probability-quantifier language \(L_{\omega P}\)
(\cref{def:p15-pq-models,def:p15-pq-language}).
\emph{Conclusions.} A placement of every semantic parameter
(\cref{def:p03-w03-p03-053-regimes}); proofs that a code-space topology
and a finite random-model probability are not measure or category
quantifiers (\cref{thm:p03-w03-p03-053-code-carrier,%
thm:p03-w03-p03-053-almost-sure,thm:p03-w03-p03-053-finite-collapse});
a measure-versus-category separation on one structure
(\cref{thm:p03-w03-p03-053-separation}); failure of ordinary isomorphism
invariance and of downward L\"owenheim--Skolem for the measure and
category quantifiers
(\cref{thm:p03-w03-p03-053-invariance,thm:p03-w03-p03-053-no-countable});
failure of compactness and of the countable-carrier
L\"owenheim--Skolem theorem for unrestricted \(L_2\)
(\cref{thm:p03-w03-p03-053-l2-noncompact,thm:p03-w03-p03-053-l2-uncountable});
failure of Craig interpolation for unrestricted \(L_2\) over all
topological structures (\cref{thm:p03-w03-p03-053-l2-interpolation});
the reduction of \(L_t\)-consequence to two-sorted first-order
consequence (\cref{thm:p03-w03-p03-053-lt-reduction}); the failure of
compactness, completeness and interpolation when the finite sample spaces
are read as a semantics (\cref{thm:p03-w03-p03-053-finite-row}); and a
table of transfer and nontransfer (\cref{sec:p03-w03-p03-053-summary}).

Barwise's introduction to the volume in which the topological, Borel and
probability-quantifier chapters appeared already separates these
settings: topological logics change the class of structures, while
Friedman's measure and category logics on the reals and Keisler's logic
of probability spaces are named as distinct programmes
\cite[\S2.8, p.~21]{Barwise1985BackgroundAims}.

\section{Where each semantic parameter lives}
\label{sec:p03-w03-p03-053-regimes}

\begin{hbdefinition}{def:p03-w03-p03-053-polish-carrier}{Polish-carrier structures and the two large-section quantifiers}
Let \(L\) be an at most countable first-order signature.
A \emph{Polish-carrier \(L\)-structure} is an \(L\)-structure
\(\mathcal M\) whose carrier \(X\) is a nonempty Polish space, whose
\(n\)-ary relations are Borel subsets of \(X^n\), and whose \(n\)-ary
functions are Borel maps \(X^n\to X\). This is the injective case of a
Borel presentation in \cite[Definition~1.1, p.~1]{MontalbanNiesBorelSurvey},
with the identity as the equivalence relation.

A \emph{category structure} is such an \(\mathcal M\), read with the
Polish topology of \(X\). A \emph{measure structure} is a pair
\((\mathcal M,\mu)\) with \(\mu\) a Borel probability measure on \(X\).
A subset of \(X\) is \emph{meager} if it is contained in a countable
union of nowhere dense sets, and \emph{comeager} if its complement is
meager. A subset is \emph{\(\mu\)-null} if it is contained in a Borel
set of \(\mu\)-measure zero, and \emph{has \(\mu\)-measure one} if its
complement is \(\mu\)-null.

The language \(L(Q_c)\), respectively \(L(Q_m)\), is first-order
\(L\) with one more formation rule: if \(\varphi\) is a formula and
\(x\) a variable, then \(Q_cx\,\varphi\), respectively
\(Q_mx\,\varphi\), is a formula in which \(x\) is bound.
For an assignment \(s\), write
\(\varphi(\mathcal M,s,x)=\{a\in X:\mathcal M\models\varphi[s(x/a)]\}\)
for the section. Then
\[
\begin{aligned}
 \mathcal M\models Q_cx\,\varphi\,[s]
   &\iff \varphi(\mathcal M,s,x)\ \text{is comeager in }X,\\
 (\mathcal M,\mu)\models Q_mx\,\varphi\,[s]
   &\iff \varphi(\mathcal M,s,x)\ \text{has \(\mu\)-measure one}.
\end{aligned}
\]
All other clauses are the first-order clauses on the carrier \(X\).
\end{hbdefinition}

These are the large-section operators of
\cite[\S2.4, pp.~7--8]{MioSkrzypczakMichalewski2018}, applied to the
satisfaction set of a formula inside one structure; the meager,
comeager and null notions are those of
\cite[\S2.2, p.~5 and \S2.3, p.~6]{MioSkrzypczakMichalewski2018}.
The clauses make sense for every section, measurable or not, as noted in
\cite[Remark~2.7, p.~8]{MioSkrzypczakMichalewski2018}.
Montalb\'an and Nies describe Friedman's quantifiers in the same
terms, as ``for all \(x\) in a co-meager set'' and the full-measure
analogue, and state that Friedman studied axiomatizations and
completeness for these languages
\cite[p.~2]{MontalbanNiesBorelSurvey}. The definition above fixes the
general Polish-carrier regime used in this chapter. It does not fix the
carrier, the class of admissible measures or the axioms of Friedman's
or Steinhorn's own measure and category logics; those belong to
Steinhorn's chapter \cite[ch.~XVI, from p.~579]{Steinhorn1985BorelMeasureCategory},
whose axiomatizability and completeness results are presented with the
Borel model theory of this chapter. The results below are proved for
the general regime and therefore hold in every special case of it.

\begin{hbdefinition}{def:p03-w03-p03-053-regimes}{Six semantic regimes}
The following regimes are compared. In each, ``parameter'' names the
additional mathematical object a truth definition consults.
\begin{enumerate}
\item[(R1)] \emph{Ordinary first-order logic} on ordinary
  \(L\)-structures. No parameter beyond the structure.
\item[(R2)] \emph{\(L_{\omega_1\omega}\)} on ordinary \(L\)-structures
  (\cref{def:p02-lomega-one-omega,def:p03-morley-infinitary}).
  No parameter; only the formulas are longer.
\item[(R3)] \emph{\(L_2\) and \(L_t\)} on topological structures
  \((X,\tau,\ldots)\)
  (\cref{def:p03-topological-carrier,def:p03-topological-languages}).
  Parameter: a topology \(\tau\) on the carrier of the one structure.
  \(L_2\) quantifies over all of \(\tau\); \(L_t\) only through the
  guarded, polarity-restricted open quantifiers. With an arbitrary point
  signature this is the setting of
  \cref{sec:p03-topological-reduced-products}.
\item[(R4)] \emph{The coding space} \(X_L\) of
  \cref{def:p03-morley-coding}. Parameter: the product topology on the
  set of codes of \(L\)-structures with domain \(\omega\). It is fixed by
  \(L\) alone. The formulas evaluated at a code are ordinary
  first-order or \(L_{\omega_1\omega}\) formulas, and
  \cref{thm:p03-morley-borel-truth} describes the set of codes that
  satisfy one of them.
\item[(R5)] \emph{\(L(Q_c)\) and \(L(Q_m)\)} on Polish-carrier
  structures (\cref{def:p03-w03-p03-053-polish-carrier}).
  Parameter: the Polish topology, respectively a Borel probability
  measure, on the carrier of the one structure. It is consulted only
  through the new quantifier, on sections of that carrier.
\item[(R6)] \emph{Finite probability.} (a) The zero--one law
  (\cref{sec:p03-zero-one-law}): ordinary first-order sentences, and the
  uniform probability on the finite set of labelled structures with
  universe \([n]\). (b) \(L_{\omega P}\)
  (\cref{def:p15-pq-models,def:p15-pq-language}): a threshold binder
  \((\exists^{\ge r}\bar y)\), \(0<r<1\), interpreted by the
  normalized counting measure on tuples from the carrier of one finite
  structure, together with a separate model measure \(\mu_n\) on the
  set of all structures of size \(n\).
\end{enumerate}
\end{hbdefinition}

\begin{center}
\small
\begin{tabular}{@{}p{0.15\linewidth}p{0.22\linewidth}p{0.24\linewidth}p{0.24\linewidth}@{}}
\hline
Regime & Parameter & Lives on & Read by \\
\hline
(R1), (R2) & none & --- & --- \\
(R3) & topology \(\tau\) & carrier of one structure & open-sort quantifiers \\
(R4) & product topology & set of codes of structures & no formula; a property of a set of codes \\
(R5) & topology or measure \(\mu\) & carrier of one structure & \(Q_c\) or \(Q_m\) on sections \\
(R6a) & uniform probability & set of labelled size-\(n\) structures & no formula; a limit of model probabilities \\
(R6b) & counting measure on tuples & carrier of one finite structure & threshold binders \((\exists^{\ge r}\bar y)\) \\
\hline
\end{tabular}
\end{center}

The table is the content of the comparison in compressed form. Two
rows that look alike, a topology in (R3) and in (R4), act on different
objects. So do the two probabilities in (R5) and (R6a).
The remaining sections prove that these differences change truth values.

\section{Codes, carriers and random finite models}
\label{sec:p03-w03-p03-053-codes-carriers}

\begin{hbtheorem}{thm:p03-w03-p03-053-code-carrier}{proved=annex:p03-w03-p03-053-code-carrier}{The code determines no carrier topology}
Let \(L\) be an at most countable signature and let \(\mathcal M\) be a
code in \(X_L\), or in \(X_{L,m}\) with \(m\ge2\). Let
\(\mathcal M^{\mathrm{dis}}\) and \(\mathcal M^{\mathrm{ind}}\) be the
topological structures with the same \(L\)-structure on the same
carrier and, respectively, the discrete and the indiscrete topology.
Then:
\begin{enumerate}
\item \(\mathcal M^{\mathrm{dis}}\) and \(\mathcal M^{\mathrm{ind}}\)
  satisfy exactly the same \(L_{\omega_1\omega}\)-sentences of \(L\),
  and both lie over the same point of \(X_L\);
\item the \(L_t\)-sentence \(T_1\) of
  \cref{thm:p03-topological-property-formulas} is true in
  \(\mathcal M^{\mathrm{dis}}\), and \(T_0\) is false in
  \(\mathcal M^{\mathrm{ind}}\).
\end{enumerate}
Hence no function of the code alone determines \(L_t\)-truth, and the
topology of \(X_L\) is not a topology on the carrier of any coded
structure.
\end{hbtheorem}

\begin{hbtheorem}{thm:p03-w03-p03-053-almost-sure}{proved=annex:p03-w03-p03-053-almost-sure}{The almost-sure theory has no finite model}
Fix the signature and the uniform labelled distribution of
\cref{thm:p03-fol-zero-one}, and write \(\mu_n(\varphi)\) for the
probability that a uniformly random structure with universe \([n]\)
satisfies \(\varphi\). Put
\[
 T_{\mathrm{as}}=\Bigl\{\varphi:\ \varphi\ \text{a first-order sentence and}
   \ \lim_{n\to\infty}\mu_n(\varphi)=1\Bigr\}.
\]
Then \(T_{\mathrm{as}}\) is complete: for every sentence \(\varphi\),
exactly one of \(\varphi,\neg\varphi\) belongs to it. Every finite
subset of \(T_{\mathrm{as}}\) is true in some finite structure;
\(T_{\mathrm{as}}\) itself is satisfiable; and \(T_{\mathrm{as}}\) has no
finite model. In particular \(T_{\mathrm{as}}\) is not the theory of
any finite structure: for every structure in any of the sample spaces
there is a sentence whose limit probability differs from its truth
value in that structure.
\end{hbtheorem}

\begin{hbtheorem}{thm:p03-w03-p03-053-finite-collapse}{proved=annex:p03-w03-p03-053-finite-collapse}{Measure and category quantifiers collapse on finite carriers}
Let \(\mathcal M\) be a Polish-carrier structure with a finite carrier
\(X\). Then the topology of \(X\) is discrete, and for every formula
\(\varphi\) and assignment \(s\),
\[
 \mathcal M\models Q_cx\,\varphi\,[s]
   \iff \mathcal M\models\forall x\,\varphi\,[s].
\]
If \(\mu\) is a probability measure on \(X\) that gives every point
positive mass, for instance the normalized counting measure, then
likewise
\[
 (\mathcal M,\mu)\models Q_mx\,\varphi\,[s]
   \iff \mathcal M\models\forall x\,\varphi\,[s].
\]
The threshold binders of \(L_{\omega P}\) are not of this form: for
every \(r\) with \(0<r<1\) there is a finite structure with a unary
relation \(R\) in which \((\exists^{\ge r}y)R(y)\) is true and
\(\forall y\,R(y)\) is false.
\end{hbtheorem}

\begin{hbwarning}{def:p03-w03-p03-053-no-substitution}{Neither code topology nor random finite models is a measure or category logic}
A measure or category logic in the sense of (R5) evaluates its new
quantifier \emph{inside one structure}, on a section of that
structure's carrier. Neither of the following does so.
\begin{enumerate}
\item The topology of \(X_L\) in \cref{sec:p03-morley-borel-truth}
  is a topology on a set of structures. The Borel truth sets of
  \cref{thm:p03-morley-borel-truth} are sets of codes; they contain
  no information about neighbourhoods of elements
  (\cref{thm:p03-w03-p03-053-code-carrier}). One may put a measure or a
  category notion on \(X_L\) and ask whether a Borel truth set is
  large; that is a statement \emph{about} a class of countable
  structures, not the truth of a formula in any of them.
\item The probability in \cref{sec:p03-zero-one-law} is a probability on
  the finite set of labelled size-\(n\) structures. The sentences are
  ordinary first-order sentences, and the limit law describes a
  complete theory with no finite model
  (\cref{thm:p03-w03-p03-053-almost-sure}).
  On the finite carriers of that sample space the measure and category
  quantifiers would collapse to \(\forall\)
  (\cref{thm:p03-w03-p03-053-finite-collapse}).
\end{enumerate}
The probability-quantifier language \(L_{\omega P}\) does read a measure
on a carrier, but it is the normalized counting measure on a finite
carrier, with thresholds strictly between zero and one; its asymptotic
statements use the separate model measure \(\mu_n\)
(\cref{def:p15-pq-models}). Unrestricted \(L_{\omega P}\) has no
zero--one law for that model measure (\cref{thm:p15-pq-limit-failures}),
so the zero--one law of (R6a) does not transfer to (R6b).
\end{hbwarning}

\section{Measure and category on one Borel structure}
\label{sec:p03-w03-p03-053-measure-category}

Write \(2^\omega\) for the Cantor space of infinite binary sequences
with the product of discrete topologies, and \([t]\) for the basic
clopen set of sequences extending a finite binary string \(t\).
Let \(\lambda\) be the coin-flipping measure, the Borel probability
measure on \(2^\omega\) under which each coordinate takes each value
with probability \(1/2\), independently; in particular
\(\lambda([t])=2^{-|t|}\). Its existence and uniqueness are taken from
\cite[\S2.3, p.~6]{MioSkrzypczakMichalewski2018}, where it is called
Lebesgue measure on \(\omega\)-words. Each unary relation below is a
Borel subset, so \((2^\omega,P)\) is a Polish-carrier structure.

\begin{hblemma}{thm:p03-w03-p03-053-baire}{proved=annex:p03-w03-p03-053-baire}{Baire category theorem}
In a nonempty complete metric space no nonempty open set is meager.
Consequently, in a nonempty Polish space the empty set is not comeager,
a countable intersection of comeager sets is comeager, and a nonempty
countable Polish space has an isolated point.
\end{hblemma}

\begin{hbtheorem}{thm:p03-w03-p03-053-separation}{proved=annex:p03-w03-p03-053-separation}{A comeager null section}
For \(n\ge1\) let
\(B_n=\{x\in2^\omega:x(k)=0\ \text{for all } k \text{ with } n\le k\le2n\}\)
and put
\[
  G=\bigcap_{m\ge1}\ \bigcup_{n\ge m}B_n .
\]
Then \(G\) is a Borel set, comeager in \(2^\omega\), and
\(\lambda(G)=0\). Hence, in the structure \((2^\omega,G)\) with unary
predicate \(G\), both
\[
  Q_cx\,G(x)\qquad\text{and}\qquad Q_mx\,\neg G(x)
\]
are true, the first for the Cantor topology and the second for
\(\lambda\). Within either quantifier alone, for all formulas
\(\varphi,\psi\), the formulas
\(Q x\,\varphi\land Qx\,\psi\to Qx(\varphi\land\psi)\) and
\(\neg(Qx\,\varphi\land Qx\,\neg\varphi)\) are true under every
assignment in every structure of regime (R5), where \(Q\) is \(Q_c\)
throughout or \(Q_m\) throughout. The mixed formula
\(Q_cx\,\varphi\to\neg Q_mx\,\neg\varphi\) is false in
\((2^\omega,G)\), read with both the Cantor topology and \(\lambda\),
for \(\varphi=G(x)\).
\end{hbtheorem}

\begin{hbtheorem}{thm:p03-w03-p03-053-invariance}{proved=annex:p03-w03-p03-053-invariance}{Ordinary isomorphism does not preserve the new quantifiers}
Let \(P_1=\{x\in2^\omega:x(0)=0\}\) and
\(P_2=\{x\in2^\omega:x(2k)=0\ \text{for all } k\}\).
The structures \((2^\omega,P_1)\) and \((2^\omega,P_2)\) are isomorphic
as ordinary \(\{P\}\)-structures, so they satisfy the same
\(L_{\omega_1\omega}\)-sentences. Yet \(Q_mx\,\neg P(x)\) and
\(Q_cx\,\neg P(x)\) are false in the first and true in the second.
Conversely, a bijection between Polish-carrier structures that is an
\(L\)-isomorphism preserves and reflects every \(L(Q_c)\)-formula if it
is a homeomorphism, and every \(L(Q_m)\)-formula if it and its inverse
are Borel and it carries the one measure to the other.
\end{hbtheorem}

The new quantifiers are therefore not generalized quantifiers of the
kind obtained by adding a class of ordinary structures, whose truth
depends only on the isomorphism type of the structure.

\section{Compactness, L\"owenheim--Skolem, completeness and interpolation}
\label{sec:p03-w03-p03-053-transfer}

\begin{hbtheorem}{thm:p03-w03-p03-053-no-countable}{proved=annex:p03-w03-p03-053-no-countable}{No countable carrier for ``every point is negligible''}
Let
\(\sigma_c=\forall y\,Q_cx\,(x\neq y)\) and
\(\sigma_m=\forall y\,Q_mx\,(x\neq y)\), in the empty signature.
In regime (R5), every model of \(\sigma_c\) and every model of
\(\sigma_m\) has an uncountable carrier, while both are true in
\(2^\omega\) with its topology and with \(\lambda\). Hence the
countable-model form of downward L\"owenheim--Skolem
(\cref{thm:downward-ls}) fails for \(L(Q_c)\) and \(L(Q_m)\) in
this regime, already for one sentence.
\end{hbtheorem}

By contrast, \(L_t\) keeps a countable form of the theorem: a satisfiable
countable-signature set of \(L_t\)-sentences has a model with countably
many points and a countable basis, although the topology generated by
that basis need not be countable
(\cref{thm:p03-topological-compactness-ls}).
Unrestricted \(L_2\) does not keep it.

\begin{hbtheorem}{thm:p03-w03-p03-053-l2-uncountable}{proved=annex:p03-w03-p03-053-l2-uncountable}{An \(L_2\)-sentence with only uncountable carriers}
Let the point signature contain a binary relation symbol \(<\), as in
\cref{sec:p03-topological-reduced-products}. Let \(\theta_{\mathbb R}\) be
the conjunction of the following \(L_2\)-sentences:
\[
\begin{aligned}
 \mathrm{Lin}&:\ \text{\(<\) is a strict linear order},\\
 \mathrm{Ne}&:\ \exists x\,(x=x),\\
 \mathrm{Dense}&:\ \forall x\forall y\bigl(x<y\to\exists z(x<z\land z<y)\bigr),\\
 \mathrm{NoEnd}&:\ \forall x\exists y\exists z(y<x\land x<z),\\
 \mathrm{Int}&:\ \forall a\forall b\,\exists U\,\forall y
     \bigl(y\varepsilon U\leftrightarrow(a<y\land y<b)\bigr),\\
 \mathrm{Sup}&:\ \forall U\Bigl(\bigl(\exists x\,x\varepsilon U\land
     \exists b\,\forall y(y\varepsilon U\to y<b)\bigr)\\
   &\qquad\to\exists s\bigl(\forall y(y\varepsilon U\to\neg\,s<y)\land
     \forall t\,(t<s\to\exists y(y\varepsilon U\land t<y))\bigr)\Bigr).
\end{aligned}
\]
Then the real line with its usual order and its usual topology is a
topological model of \(\theta_{\mathbb R}\), and every topological model
of \(\theta_{\mathbb R}\) has an uncountable carrier. Hence the
countable-carrier form of downward L\"owenheim--Skolem fails for
\(L_2\) on topological structures, already for one sentence.
\end{hbtheorem}

The sentence \(\mathrm{Int}\) quantifies existentially over open sets
without a membership guard, and \(\mathrm{Sup}\) quantifies universally
over all open sets; neither is an \(L_t\)-sentence, as the countable form
of the theorem for \(L_t\) requires.

\begin{hbtheorem}{thm:p03-w03-p03-053-l2-noncompact}{proved=annex:p03-w03-p03-053-l2-noncompact}{Unrestricted \(L_2\) is not compact on topological structures}
Let the point signature contain a binary relation symbol \(<\), as in
\cref{sec:p03-topological-reduced-products}. Let \(\mathrm{Lin}\) say
that \(<\) is a strict linear order, let \(\mathrm{Discrete}\) be the
\(L_t\)-sentence of \cref{thm:p03-topological-property-formulas}, let
\(\exists_{\ge n}\) say that there are at least \(n\) points, and let
\[
\begin{aligned}
 \mathrm{Max}(U)&=\exists x\bigl(x\varepsilon U\land
                   \forall y(y\varepsilon U\to\neg x<y)\bigr),\\
 \mathrm{Min}(U)&=\exists x\bigl(x\varepsilon U\land
                   \forall y(y\varepsilon U\to\neg y<x)\bigr),\\
 \mathrm{Fin}&=\forall U\bigl(\exists x\,x\varepsilon U\to
                   (\mathrm{Max}(U)\land\mathrm{Min}(U))\bigr).
\end{aligned}
\]
Then the set of \(L_2\)-sentences
\[
 \Gamma=\{\mathrm{Lin},\mathrm{Discrete},\mathrm{Fin}\}\cup
        \{\exists_{\ge n}:n\ge1\}
\]
is finitely satisfiable in topological structures but has no
topological model.
Consequently there is no derivability relation \(\vdash\) between sets
of \(L_2\)-sentences and \(L_2\)-sentences that is sound and strongly
complete for consequence over topological structures and has finite
character: \(\Delta\vdash\psi\) implies \(\Delta_0\vdash\psi\) for some
finite \(\Delta_0\subseteq\Delta\).
\end{hbtheorem}

The sentence \(\mathrm{Fin}\) uses an unguarded open quantifier, so it is
not an \(L_t\)-sentence; \(L_t\) is compact
(\cref{thm:p03-topological-compactness-ls}). In \(L_{\omega_1\omega}\)
compactness fails for a different reason, the countable disjunction
(\cref{thm:p02-lomega-one-omega-calibration}), and the completeness
theorem that survives there uses infinitary derivations
(\cref{thm:p02-karp-completeness}).

\begin{hbtheorem}{thm:p03-w03-p03-053-lt-reduction}{proved=annex:p03-w03-p03-053-lt-reduction}{Reduction of \(L_t\)-consequence}
Let \(\mathrm{Ax}\) be the set of the four basis axioms Ext, Union,
Empty and Bas of \cref{def:p03-topological-basis-axioms}. For
every set \(\Gamma\cup\{\varphi\}\) of \(L_t\)-sentences,
\(\varphi\) is true in every topological model of \(\Gamma\) if and only
if \(\varphi\) is a two-sorted first-order consequence of
\(\Gamma\cup\mathrm{Ax}\). Consequently \(L_t\)-consequence over
topological structures is compact, and it is captured by any sound and
strongly complete calculus for two-sorted first-order logic with the
four axioms added. For a recursive signature the set of \(L_t\)-sentences
true in all topological structures is recursively enumerable.
\end{hbtheorem}

For unrestricted \(L_2\), \cref{thm:p03-w03-p03-053-l2-noncompact} shows
that no such reduction to a finitary first-order consequence relation
exists.

\begin{hbexample}{def:p03-w03-p03-053-cantor-example}{One space read as codes and as a carrier}
Let \(L=\{P\}\) with \(P\) unary. By \cref{def:p03-morley-coding},
\(X_L=2^\omega\): a point \(x\) codes the structure
\((\omega,\{k:x(k)=1\})\).
\begin{enumerate}
\item \emph{Code reading (R4).} The first-order sentence
  \(\exists x\,P(x)\) has truth set \(2^\omega\setminus\{0^\omega\}\),
  an open set of codes. It is comeager and has \(\lambda\)-measure one.
  This says that almost every countable \(\{P\}\)-structure, in either
  sense of ``almost'', has a nonempty predicate. It is a statement about
  a set of structures; each coded structure satisfies or fails
  \(\exists x\,P(x)\) outright. The set \(G\) of
  \cref{thm:p03-w03-p03-053-separation} is, in this reading, a Borel set
  of codes that is not closed under isomorphism of the coded structures,
  as shown after this example.
\item \emph{Carrier reading (R5).} The same \(G\) is the predicate of
  the single structure \((2^\omega,G)\). Its elements are the points of
  \(2^\omega\); \(Q_cx\,G(x)\) is a true \(L(Q_c)\)-sentence about it
  for the Cantor topology, and \(Q_mx\,\neg G(x)\) is a true
  \(L(Q_m)\)-sentence about it for \(\lambda\).
\item \emph{Carrier topology (R3).} With the Cantor topology,
  \((2^\omega,G)\) is a topological structure satisfying the
  \(L_t\)-sentence \(T_2\), because two sequences that first differ at
  coordinate \(k\) lie in disjoint basic sets \([t]\) with \(|t|=k+1\).
  With the indiscrete topology on the same carrier it fails \(T_0\), as
  in \cref{thm:p03-w03-p03-053-code-carrier}.
\end{enumerate}
\end{hbexample}

The isomorphism remark in the code reading is justified as follows. The
code \(x\) of a structure with infinite coinfinite predicate lies in
\(G\) exactly when \(x\) vanishes on the whole run \(x(n),\ldots,x(2n)\) for
infinitely many \(n\). Arrange the predicate on the even numbers: the
resulting code has no such run for \(n\ge1\), so it lies outside
\(G\). Arrange it on the set of numbers \(j\) such that \(j=2^{2i}\) for some \(i\):
for each \(i\ge1\) the number \(n=2^{2i}+1\) satisfies
\(2^{2i}<n\) and \(2n<2^{2i+2}\), so the run from \(n\) to \(2n\)
contains no power of four and is zero. There are infinitely many such
\(n\), so
that code lies in \(G\). Both codes describe structures with countably
infinite predicate and countably infinite complement, hence isomorphic
ones.

\section{The finite sample spaces as a semantics}
\label{sec:p03-w03-p03-053-finite}

The structures in the sample spaces of (R6) are finite. If one reads
first-order logic, or \(L_{\omega P}\), with consequence restricted to
nonempty finite structures, the ordinary theorems change as follows.
Compactness fails already for first-order sentences
(\cref{def:p03-finite-compactness-failure}), and finite validity is not
computably enumerable (\cref{thm:trakhtenbrot}). Interpolation fails as
well. The proof uses the following lemma on the empty signature.

\begin{hblemma}{thm:p03-w03-p03-053-pure-equality}{proved=annex:p03-w03-p03-053-pure-equality}{Large finite sets agree on sentences of small rank}
Let \(q\ge0\), and let \(A\) and \(B\) be finite sets with
\(|A|\ge q\) and \(|B|\ge q\). Then \(A\) and \(B\), as structures of
the empty signature, satisfy the same first-order sentences of
quantifier rank at most \(q\), all of whose atoms are equations between
variables.
\end{hblemma}

\begin{hbtheorem}{thm:p03-w03-p03-053-finite-row}{proved=annex:p03-w03-p03-053-finite-row}{First-order and probability-quantifier logic over finite structures}
Write \(\Gamma\models_{\mathrm{fin}}\psi\) when \(\psi\) is true in every
nonempty finite structure, for the joint signature, in which every
member of \(\Gamma\) is true.
\begin{enumerate}
\item Compactness fails for \(\models_{\mathrm{fin}}\), both for
  first-order logic and for \(L_{\omega P}\).
\item For every finite relational signature containing a binary relation
  symbol, the set of first-order sentences \(\psi\) with
  \(\varnothing\models_{\mathrm{fin}}\psi\) is not computably enumerable,
  and neither is the set of such \(L_{\omega P}\)-sentences whose
  thresholds are rational.
\item Craig interpolation fails for first-order logic under
  \(\models_{\mathrm{fin}}\). Let \(\mathrm{succ}(x,y)\) abbreviate
  \(x<y\land\neg\exists z(x<z\land z<y)\), and let
  \[
  \begin{aligned}
   \varphi_1={}&\mathrm{Lin}(<)\land
     \forall x\bigl(\neg\exists y\,(y<x)\to P(x)\bigr)\land
     \forall x\forall y\bigl(\mathrm{succ}(x,y)\to
       (P(x)\leftrightarrow\neg P(y))\bigr)\\
     &\land\forall x\bigl(\neg\exists y\,(x<y)\to P(x)\bigr),
  \end{aligned}
  \]
  and let \(\varphi_2\) be the same sentence in the symbols \(<',P'\)
  with the last conjunct replaced by
  \(\forall x\bigl(\neg\exists y\,(x<'y)\to\neg P'(x)\bigr)\).
  Then \(\varphi_1\models_{\mathrm{fin}}\neg\varphi_2\), and there is no
  first-order sentence \(\chi\) of the empty signature, the only
  signature \(\{<,P\}\) and \(\{<',P'\}\) share, with
  \(\varphi_1\models_{\mathrm{fin}}\chi\) and
  \(\chi\models_{\mathrm{fin}}\neg\varphi_2\).
\end{enumerate}
\end{hbtheorem}

In a nonempty finite linear order the sentence \(\varphi_1\) forces
\(P\) onto the elements in odd positions and then says that the last
position is odd; \(\varphi_2\) says that the size is even. A sentence of
pure equality cannot tell large odd sizes from large even ones
(\cref{thm:p03-w03-p03-053-pure-equality}). Over all structures the same
pair has no such consequence: the integers, with \(<\) and \(<'\) both the
usual order and \(P\) and \(P'\) both the set of even integers, satisfy
\(\varphi_1\land\varphi_2\), because there is no first and no last element
and consecutive integers differ in parity. So the example does not
conflict with first-order Craig interpolation
(\cref{thm:p02-craig-interpolation}), which concerns \(\models\), not
\(\models_{\mathrm{fin}}\).

\section{Interpolation for unrestricted \texorpdfstring{\(L_2\)}{L2}}
\label{sec:p03-w03-p03-053-l2-interpolation}

Unrestricted \(L_2\) does not need a restriction to finite structures to
lose interpolation: its unguarded open quantifiers can force a carrier to
be finite. The pair of sentences used for \(\models_{\mathrm{fin}}\) then
works for consequence over all topological structures. The common part of
the two sentences is the topology alone, and on discrete spaces an
\(L_2\)-sentence with no point symbols cannot distinguish two large finite
sizes.

Throughout this section the point signature may contain relation symbols
of any arity, as in \cref{sec:p03-topological-reduced-products}. For a
point signature \(L\), write \(L_2(L)\) for the \(L_2\)-formulas whose
point symbols lie in \(L\); its atoms are point equality, open equality,
membership \(x\varepsilon U\) and the atoms of \(L\). The quantifier rank
of an \(L_2\)-formula counts nested quantifiers of both sorts. A
topological structure is \emph{discrete} if its topology is the power set
of its carrier.

\begin{hbdefinition}{def:p03-w03-p03-053-l2-craig}{Craig interpolation for \texorpdfstring{\(L_2\)}{L2}}
For point signatures \(L^1\) and \(L^2\), write
\(\varphi\models_{\mathrm{top}}\psi\) when \(\psi\) is true in every
topological \((L^1\cup L^2)\)-structure in which \(\varphi\) is true.
\(L_2\) has the \emph{Craig interpolation property} over topological
structures if for all such \(L^1,L^2\) and all \(\varphi\in L_2(L^1)\),
\(\psi\in L_2(L^2)\) with \(\varphi\models_{\mathrm{top}}\psi\) there is
\(\chi\in L_2(L^1\cap L^2)\) with \(\varphi\models_{\mathrm{top}}\chi\) and
\(\chi\models_{\mathrm{top}}\psi\). The topology is part of every
structure, so open quantifiers, open equality and membership are
available to \(\chi\) whatever \(L^1\cap L^2\) is.
\end{hbdefinition}

\begin{hblemma}{thm:p03-w03-p03-053-discrete-rank}{proved=annex:p03-w03-p03-053-discrete-rank}{Large finite discrete spaces agree on sentences of small rank}
Let \(q\ge0\), and let \((X,\mathcal P(X))\) and \((Y,\mathcal P(Y))\) be
finite discrete topological structures for the empty point signature with
\(|X|\ge2^q\) and \(|Y|\ge2^q\). Then they satisfy the same sentences of
\(L_2(\varnothing)\) of quantifier rank at most \(q\).
\end{hblemma}

In \cref{thm:p03-w03-p03-053-pure-equality} the threshold is \(q\); here
it is \(2^q\), because an open quantifier may cut a block of points into
two parts and each part must still be large for the remaining
quantifiers.

\begin{hbtheorem}{thm:p03-w03-p03-053-l2-interpolation}{proved=annex:p03-w03-p03-053-l2-interpolation}{Craig interpolation fails for unrestricted \texorpdfstring{\(L_2\)}{L2}}
Let \(\varphi_1\) and \(\varphi_2\) be the first-order sentences of
\cref{thm:p03-w03-p03-053-finite-row}(3), in the point signatures
\(L^1=\{<,P\}\) and \(L^2=\{<',P'\}\). Let \(\mathrm{Discrete}\) be the
\(L_t\)-sentence of \cref{thm:p03-topological-property-formulas}, let
\(\mathrm{Fin}\) be the \(L_2\)-sentence of
\cref{thm:p03-w03-p03-053-l2-noncompact}, and let \(\mathrm{Fin}'\) be
\(\mathrm{Fin}\) with \(<'\) in place of \(<\). Put
\[
 \psi_1=\mathrm{Discrete}\land\mathrm{Fin}\land\varphi_1,
 \qquad
 \psi_2=\mathrm{Discrete}\land\mathrm{Fin}'\land\varphi_2 .
\]
\begin{enumerate}
\item A topological \(L^1\)-structure satisfies \(\psi_1\) if and only if
  it is discrete, its carrier is finite of odd size, and \(<\) and \(P\)
  are as described for \(\varphi_1\) in
  \cref{thm:p03-w03-p03-053-finite-row}(3). Every finite discrete space
  of odd size has such an expansion. The same holds for \(\psi_2\) with
  even size, \(<'\) and \(P'\).
\item \(\psi_1\models_{\mathrm{top}}\neg\psi_2\).
\item There is no \(\chi\in L_2(\varnothing)\) with
  \(\psi_1\models_{\mathrm{top}}\chi\) and
  \(\chi\models_{\mathrm{top}}\neg\psi_2\). Hence \(L_2\) does not have
  the Craig interpolation property over topological structures.
\end{enumerate}
\end{hbtheorem}

\begin{hbexample}{def:p03-w03-p03-053-l2-interpolation-example}{Why no sentence about the topology separates the two sides}
Take \(q=2\), so \(2^q=4\). The discrete spaces with \(5\) and \(6\)
points satisfy the same \(L_2(\varnothing)\)-sentences of rank at most
\(2\) by \cref{thm:p03-w03-p03-053-discrete-rank}. The first expands to a
model of \(\psi_1\): order the points as \(e_1<\cdots<e_5\) and put
\(P=\{e_1,e_3,e_5\}\). The second expands to a model of \(\psi_2\):
order them as \(e_1<'\cdots<'e_6\) and put \(P'=\{e_1,e_3,e_5\}\), so
that \(e_6\notin P'\). A candidate interpolant of rank \(2\), being true in
the first space, is true in the second, which is a model of \(\psi_2\);
so it fails \(\chi\models_{\mathrm{top}}\neg\psi_2\). The same comparison
between an odd size \(m\ge2^q\) and \(m+1\) rules out every rank \(q\).
\end{hbexample}

The failure has two ingredients that \(L_t\) lacks. The sentence
\(\mathrm{Fin}\) quantifies universally over all open sets without a
membership guard, so it is not an \(L_t\)-sentence; and on discrete spaces
it forces finiteness, so that the parity information carried by
\(\varphi_1\) and \(\varphi_2\) concerns finite sizes, which
\(L_2(\varnothing)\) cannot separate. The theorem therefore says nothing
about interpolation for \(L_t\), which is Ziegler's subject
\cite[ch.~XV, \S2, from p.~560]{Ziegler1985TopologicalModelTheory}. Over
ordinary structures the pair \(\varphi_1,\varphi_2\) has the joint model
on the integers described above, and first-order interpolation is
untouched (\cref{thm:p02-craig-interpolation}). The counterexample uses a
binary point relation in each signature; when both point signatures
consist of unary predicates only, as in the correspondence signature of
\cref{def:p03-topological-languages}, the argument does not apply, and
no proof for that case is given here.

\section{Summary of transfer and nontransfer}
\label{sec:p03-w03-p03-053-summary}
The table records, for each logic and regime that defines a consequence
relation, what becomes of the four ordinary first-order theorems. Regimes
(R4) and (R6a) define no consequence relation of their own: in (R4) the
formulas are ordinary and the topology classifies sets of codes
(\cref{thm:p03-morley-borel-truth}); in (R6a) the limit law produces one
complete theory with no finite model
(\cref{thm:p03-w03-p03-053-almost-sure}). Their finite sample spaces, read
as a semantics, give the last row.

\begin{center}
\footnotesize
\begin{tabular}{@{}p{0.15\linewidth}p{0.16\linewidth}p{0.2\linewidth}p{0.21\linewidth}p{0.18\linewidth}@{}}
\hline
Logic, regime & Compactness & Countable models & Completeness & Interpolation \\
\hline
First-order (R1) & holds (\cref{thm:fol-compactness}) &
  holds (\cref{thm:downward-ls}) &
  strong, finitary (\cref{thm:fol-completeness}) &
  Craig holds (\cref{thm:p02-craig-interpolation}) \\
\(L_{\omega_1\omega}\) (R2) & fails
  (\cref{thm:p02-lomega-one-omega-calibration}) &
  holds for countably many sentences of a countable signature
  (\cref{tab:p02-fol-extensions}) &
  infinitary derivations (\cref{thm:p02-karp-completeness}) &
  Craig and Beth hold \cite[\S1.6, p.~10]{Barwise1985BackgroundAims} \\
\(L_t\) (R3) & holds (\cref{thm:p03-topological-compactness-ls}) &
  countable carrier and basis (\cref{thm:p03-topological-compactness-ls}) &
  via two-sorted first-order logic (\cref{thm:p03-w03-p03-053-lt-reduction}) &
  Ziegler's theorem; see below \\
\(L_2\) (R3) & fails (\cref{thm:p03-w03-p03-053-l2-noncompact}) &
  fails (\cref{thm:p03-w03-p03-053-l2-uncountable}) &
  no strongly complete finitary calculus (\cref{thm:p03-w03-p03-053-l2-noncompact}) &
  Craig fails (\cref{thm:p03-w03-p03-053-l2-interpolation}) \\
\(L(Q_c)\), \(L(Q_m)\) (R5) & Steinhorn's model classes; see below &
  fails (\cref{thm:p03-w03-p03-053-no-countable}) &
  Steinhorn's model classes; see below &
  not decided in this chapter \\
First-order and \(L_{\omega P}\) over finite structures (R6) &
  fails (\cref{thm:p03-w03-p03-053-finite-row}) &
  every model is finite &
  finite validity not computably enumerable
  (\cref{thm:p03-w03-p03-053-finite-row}) &
  Craig fails for first-order sentences
  (\cref{thm:p03-w03-p03-053-finite-row}) \\
\hline
\end{tabular}
\end{center}

The changes are of three kinds. In (R2) and (R3) the class of structures
is unchanged or enlarged by a topology, and the theorems fail through
expressive power: countable disjunctions in \(L_{\omega_1\omega}\),
unguarded quantification over all open sets in \(L_2\), which also
destroys Craig interpolation by forcing finite carriers
(\cref{thm:p03-w03-p03-053-l2-interpolation}). Restricting the
open quantifiers to the guarded form of \(L_t\) restores compactness, the
countable form of L\"owenheim--Skolem and a finitary completeness theorem,
because \(L_t\)-consequence is first-order consequence from four axioms
(\cref{thm:p03-w03-p03-053-lt-reduction}). In (R5) the carrier must carry
a Polish topology or a measure, and countable models are lost because
the new quantifiers make singletons negligible
(\cref{thm:p03-w03-p03-053-no-countable}); isomorphism invariance is lost
as well (\cref{thm:p03-w03-p03-053-invariance}). In (R6) nothing is added
to the syntax of first-order logic; the class of structures shrinks to the
finite ones, and compactness, completeness and interpolation fail for that
reason alone.

Two rows depend on sources presented elsewhere in this part. Interpolation
for \(L_t\) on topological structures is the subject of Ziegler's chapter
\cite[ch.~XV, \S2, from p.~560]{Ziegler1985TopologicalModelTheory}.
Axiomatizability, its consequences for the category and measure logics,
and completeness theorems for those logics over the Borel model classes
fixed in Steinhorn's chapter are the subjects of its second and third
sections \cite[ch.~XVI, \S2, from p.~586 and \S3, from
p.~591]{Steinhorn1985BorelMeasureCategory}. Those results are stated for
Steinhorn's own model classes and axioms, which this chapter does not fix;
the results proved here for (R5) hold for every Polish-carrier structure
and therefore in each of those classes. The entry ``not decided in this
chapter'' in the (R5) row marks an interpolation question for which no
proof is given here.
Craig interpolation for first-order logic and for \(L_{\omega_1\omega}\) is
recorded by Barwise, who credits the \(L_{\omega_1\omega}\) analogues of the
Craig and Beth theorems to Lopez-Escobar
\cite[\S1.6, p.~10]{Barwise1985BackgroundAims}.

The completeness entries also differ in kind. Barwise frames the
completeness problem for a logic as the question whether its valid
sentences are recursively enumerable, and notes that infinitary syntax
forces a generalized formulation \cite[\S1.3, pp.~6--7 and \S1.6,
p.~10]{Barwise1985BackgroundAims}. \Cref{thm:p03-w03-p03-053-lt-reduction}
is completeness in the first sense, and the finite-structure row is its
failure in the same sense; Karp's theorem is completeness in the
generalized sense.

\section{Sources}
\label{sec:p03-w03-p03-053-sources}

The regime placements and all theorems of this chapter are supplied
editorially, with complete proofs in the annex. Definitions of meager,
null, the coin-flipping measure and the large-section operators follow
the published version of Mio, Skrzypczak and Michalewski
\cite{MioSkrzypczakMichalewski2018}. The description of Friedman's
quantifiers and of Borel presentations follows the copy on Nies's web page of
the Montalb\'an--Nies survey \cite{MontalbanNiesBorelSurvey}, whose
page numbers are those of that copy, not of the published collection.
The historical placement of the topological, Borel and
probability-quantifier programmes follows Barwise's introduction
\cite{Barwise1985BackgroundAims}. Ziegler's and Steinhorn's chapters
are cited only for the location of their topics by the contents of the
volume.
